%% 02-the-algorithm.tex — How bisection works in plain language

\section{How the Algorithm Works}

\subsection*{The Big Picture}

Imagine you need to divide a pizza into eight equal slices.
You make one straight cut through the middle---now you have two halves.
Then you cut each half in half, and each quarter in half.
Eight slices, all roughly equal.

Redistricting works the same way, except instead of a pizza you have
a state, and instead of weight you are equalizing \emph{population}.

\subsection*{Step by Step}

\begin{enumerate}
  \item \textbf{Build a map.}
    The algorithm treats the state as a network of census tracts---about
    4,000 residents each---connected wherever two tracts share a border.

  \item \textbf{Make the first cut.}
    A mathematical solver (METIS, the same tool used to distribute
    workloads across supercomputers) splits the network into two halves
    of equal population.  It tries to make the cut as short as possible,
    which keeps districts compact.

  \item \textbf{Repeat until done.}
    Each half is split again, and again, until the state has exactly as
    many pieces as it has congressional seats.

  \item \textbf{No political data---ever.}
    At no point does the algorithm see voter registration, election
    results, or the home addresses of incumbents.  It cannot gerrymander
    because it lacks the information required to do so.
\end{enumerate}

\subsection*{What Goes In / What Comes Out}

\begin{table}[h]
\centering
\caption{Inputs and outputs of the redistricting algorithm.}
\label{tab:inout}
\begin{tabular}{@{}p{0.44\textwidth} p{0.44\textwidth}@{}}
\toprule
\textbf{What goes in} & \textbf{What comes out} \\
\midrule
Census tract boundaries (geography) &
  435 congressional districts \\
Population count per tract &
  Population balance within $\pm$0.5\% \\
Shared-border lengths (geometry) &
  District assignment for every tract \\
Number of seats per state &
  Maps, compactness scores, CSV data \\
\midrule
\emph{Nothing political} &
  \emph{Nothing political} \\
\bottomrule
\end{tabular}
\end{table}

\subsection*{Why Short Borders Mean Compact Districts}

The algorithm is told to make cuts as short (cheap) as possible.
Short cuts follow natural features---rivers, ridgelines, county lines---
rather than zigzagging through neighborhoods.
The result is districts that look like reasonable geographic units,
not the salamander-shaped creatures that gave gerrymandering its name.

\subsection*{Figures from the Pipeline}

The pipeline produces a round-by-round image for every state showing
how the bisection tree unfolds.
Figure~\ref{fig:mn} shows three rounds for Minnesota (eight districts):
each round halves the remaining pieces.

\begin{figure}[h]
\centering
\fbox{\parbox{0.9\textwidth}{\centering\vspace{1.2cm}
\textit{[Minnesota bisection rounds 1--3 would appear here.}\\
\textit{See \texttt{docs/figures/minnesota\_round\_\{1,2,3\}.png}}
\textit{generated by the pipeline.]}
\vspace{1.2cm}}}
\caption{Minnesota (8 districts): three rounds of recursive bisection.
Each round divides existing pieces into two equal-population halves.}
\label{fig:mn}
\end{figure}

\noindent The same process applied to Alabama (seven districts) is shown in
\texttt{docs/figures/alabama\_round\_\{1,2,3\}.png}.
